\documentclass[11pt]{article}
\usepackage[letterpaper,margin=0.9in]{geometry}
\usepackage{amsmath,amssymb,amsthm,mathtools,array,booktabs}
\usepackage[colorlinks=true,allcolors=blue]{hyperref}
\newcommand{\Tr}{\operatorname{Tr}}
\newcommand{\CS}{\mathrm{CS}}
\newcommand{\EC}{\mathrm{EC}}
\newtheorem{theorem}{Theorem}
\newtheorem{lemma}{Lemma}
\setcounter{MaxMatrixCols}{11}
\newcolumntype{L}[1]{>{\raggedright\arraybackslash}p{#1}}
\renewcommand{\theequation}{S\arabic{equation}}
\renewcommand{\thesection}{S\arabic{section}}
\renewcommand{\thetable}{S\arabic{table}}
\title{Supplemental Material:\\Extensible Causality Does Not Guarantee Pure Quantum Dilations}
\author{Chengxi Li\textsuperscript{1,2}, Hao Zhu\textsuperscript{1,2}, and
Wanzi Sun\textsuperscript{2,3}\thanks{Corresponding author. \href{mailto:wzsun@jnxy.edu.cn}{wzsun@jnxy.edu.cn}}\\[0.5em]
\small\textsuperscript{1}Beijing Institute of Petrochemical Technology, Beijing 102617, China\\
\small\textsuperscript{2}Beijing National Laboratory for Condensed Matter Physics,\\
\small Institute of Physics, Chinese Academy of Sciences, Beijing 100190, China\\
\small\textsuperscript{3}Jining University, Qufu, Shandong 273155, China}
\date{October 10, 2026; working draft}
\begin{document}
\maketitle

The main text distinguishes causal laboratory records from a pure quantum explanation of the process connecting the laboratories. This supplement gives the explicit local response protocols and the complete proof: a process witness excludes every finite pure dilation, while a common order weight and an arbitrary-dimensional positive operator identity establish causality for every finite experiment.

\section{Definitions and conventions}
The original ports are qubits in the order $A_I,A_O,B_I,B_O$: Alice's input and output, then Bob's input and output. The subscripts $I,O$ mean input and output. The Pauli matrices are $\sigma_x,\sigma_y,\sigma_z$ in the computational basis. A superscript on a Pauli matrix identifies its port; omitted ports carry the identity. Products on distinct ports denote tensor products in the stated order. Without port labels, the tensor factors are specified locally. $I$ without a subscript denotes the identity on the specified factors, and $I_T$ the identity on a space $T$. Input ancillas $A'$ and $B'$ can have arbitrary finite dimensions. Alice chooses a setting $x$ and records an outcome $a$; Bob chooses $y$ and records $b$. A process matrix $W$ assigns probabilities
\begin{equation}
 p(a,b|x,y)=\Tr\!\left[(W\otimes\Omega^{A'B'})
 (M_{a|x}^{A_IA'A_O}\otimes N_{b|y}^{B_IB'B_O})\right],
 \label{eq:born}
\end{equation}
where $p(a,b|x,y)$ is the joint conditional probability, $\Omega^{A'B'}$ the shared input density operator, and $M_{a|x},N_{b|y}$ the Choi matrices of the local operations for the indicated settings and outcomes. Superscripts on these matrices list the spaces on which they act; factors are reordered as needed. $\Tr$ is the trace, $\otimes$ the tensor product, and $\succeq0$ positive semidefiniteness; $U\preceq V$ means $V-U\succeq0$. Thus $\Omega\succeq0$ and $\Tr\Omega=1$. Each event is completely positive (CP), and summing an instrument over outcomes gives a trace-preserving (TP) channel.

The Choi convention matters for complex operations. For a CP map $\mathcal E$, define its standard Choi matrix $J_{\mathcal E}$ and the matrix $M_{\mathcal E}$ used in Eq.~\eqref{eq:born} by
\begin{equation}
 J_{\mathcal E}=\sum_{i,j}|i\rangle\langle j|\otimes
 \mathcal E(|i\rangle\langle j|),\qquad M_{\mathcal E}=J_{\mathcal E}^{T}.
 \label{eq:choi-convention}
\end{equation}
Here $|i\rangle,|j\rangle$ are input computational-basis vectors and the sum runs over that basis. The superscript $T$ is the full transpose in the input--output product basis, $*$ denotes entrywise complex conjugation, and $\dagger$ below denotes the Hermitian adjoint. Equivalently, $M_{\mathcal E}=J_{\mathcal E}^*$, since $J_{\mathcal E}$ is Hermitian. Thus $M_{\mathcal E}\succeq0$ and $\Tr_O M_{\mathcal E}=I_I$ for a channel, where $\Tr_O$ traces out the output and $I_I$ is the input identity. Equation~\eqref{eq:born} contains no additional transpose.

For an operator $Q$ and a tensor factor $T$ of dimension $d_T$, the replacement map \({}_T Q\) traces out $T$ and reinserts the normalized identity $I_T/d_T$. Multiple labels indicate replacement of all listed factors. With trivial original global past and future, the bipartite process constraints~\cite{OCB,EC} are positivity, $\Tr W=d_{A_O}d_{B_O}=4$, and
\begin{align}
 {}_{B_IB_O}W&={}_{A_OB_IB_O}W,&
 {}_{A_IA_O}W&={}_{A_IA_OB_O}W,\notag\\
 W&={}_{A_O}W+{}_{B_O}W-{}_{A_OB_O}W.
 \label{eq:process-subspace}
\end{align}
Here $A\prec B$ denotes Alice before Bob and $B\prec A$ the reverse order. An $A\prec B$ process additionally obeys \({}_{B_O}W=W\); a $B\prec A$ process obeys \({}_{A_O}W=W\). Causal separability (CS) means a sum of positive, possibly subnormalized, processes of these orders.

A normalized conditional distribution is causal if
\begin{equation}
 p=q\,p^{A\prec B}+(1-q)\,p^{B\prec A},\qquad 0\le q\le1,
 \label{eq:causal-table}
\end{equation}
where $p$ abbreviates the whole conditional table, the two superscripted tables are normalized distributions for the indicated orders, and $q$ is the probability of the Alice-first branch. It is independent of $x,y$. Subscripts $A,B$ on a distribution denote its Alice and Bob marginals: $p_A^{A\prec B}$ is independent of $y$, and $p_B^{B\prec A}$ independent of $x$. A process is extensibly causal (EC) if this holds for every finite shared input state and every finite collection of local instruments, settings, and outcomes~\cite{EC}. The decomposition may depend on the experiment. Extensible causal separability is a stronger operator property~\cite{Wechs,Wei}; it is not the premise of our theorem.

Purifiability here has the general higher-order meaning~\cite{purif}: the original process is obtained from a pure process with an added finite global past $P$ and future $F$, by inserting a fixed pure state in $P$ and tracing $F$. Their dimensions are unrestricted and are not identified with $\operatorname{rank}W$. Here a pure process preserves unitarity: inserting local unitary operations, including their local ancillary extensions, yields a global unitary transformation.

\section{Proof map and the roles of exact verification}
\label{sec:proof-map}
The universal statement is about every complete finite experiment, not about a list of tested causal games. Table~\ref{tab:proof-map} identifies the reductions that supply its physical coverage. Algebraic verification checks the final identities and positive weights; the reductions below prove that the identities apply to the original experiments. Finite complex controls detect convention errors but do not establish the universal quantifiers.

\begin{table}[ht]
\centering\small
\begin{tabular}{@{}L{0.21\linewidth}L{0.27\linewidth}L{0.43\linewidth}@{}}
\toprule
Experimental object & Representation & Preserved quantities and scope \\
\midrule
Complex CP instruments & Full-transpose Choi convention and the positive conditional output & Original Born probabilities and the $\sigma_y$ sign; Eqs.~\eqref{eq:choi-convention}, \eqref{eq:sigma-choi}, and \eqref{eq:kernel}. \\
Input and output flags & CP refinement of the reports & Normalization and the original table after coarse graining; no assumption of diagonal original instruments. \\
Multiple Kraus branches & Rectangular Stinespring amplitudes with retained environment & Mixed event outputs and all quadratic pairings; zero padding adds no orthogonality relation; Eq.~\eqref{eq:Alice-contraction}. \\
General Bob effects & Fixed local Naimark input embedding & The original input qubit, every Pauli pairing, the joint probability, and the Hall cost; Eq.~\eqref{eq:Naimark}. \\
All setting pairs & Common first-party marginals and a fixed order weight & A causal model for the complete experiment, including every outcome subset; Eq.~\eqref{eq:prescribed-margins} and the transportation lemma. \\
Arbitrary finite auxiliary spaces & An operator sum of squares (SOS) modulo projector identities & Positivity in every finite representation; only positive Alice slack, Bob projector relations, and cross-laboratory commutation; Eq.~\eqref{eq:SOS-identity}. \\
Mixed shared states & Convex combination or a local purification & A common order weight and arbitrary finite ancillas; Section~\ref{sec:universal}. \\
\bottomrule
\end{tabular}
\caption{Physical coverage of the dimension-independent proof. The original laboratory ports remain qubits throughout.}
\label{tab:proof-map}
\end{table}

\section{Physicality and an exact witness}
The process in the main text is
\begin{equation}
 W_* =\frac14\left[I+\frac15\sigma_z^{A_I}
 +\frac3{10}\sigma_z^{A_I}\sigma_x^{B_I}\sigma_z^{B_O}
 +\frac15\sum_{k=x,y,z}\sigma_k^{A_O}\sigma_k^{B_I}\right].
 \label{eq:Wstar}
\end{equation}
Here $k$ labels the three Pauli axes, and the asterisk labels the chosen counterexample. Every nonidentity term lies in the linear subspace~\eqref{eq:process-subspace}; all are traceless, so $\Tr W_*=4$. To check positivity, first fix the eigenvalues $s=\pm1$ of $\sigma_z^{A_I}$ and $j=\pm1$ of $\sigma_z^{B_O}$. On $A_OB_I$, $\sigma_x\otimes\sigma_x$ is conserved. In its sector $t=\pm1$, the remaining eigenvalues are
\begin{equation}
 \lambda_{s,t,\pm}=
 \frac{1+s\delta+t\beta\pm
 \sqrt{\alpha^2+\beta^2(1-t)^2}}4,
 \quad (\delta,\alpha,\beta)=\left(\frac15,\frac3{10},\frac15\right).
 \label{eq:spectrum-formula}
\end{equation}
Here $\delta,\alpha,\beta$ are the coefficients of $\sigma_z^{A_I}$, $\sigma_z^{A_I}\sigma_x^{B_I}\sigma_z^{B_O}$, and the three-term Pauli sum, respectively; $\lambda_{s,t,\pm}$ labels the two eigenvalue branches in each sector. They do not depend on $j$. The full spectrum, with multiplicities in parentheses, is
\begin{equation}
 \frac1{40}(2),\quad\frac18(2),\quad\frac7{40}(2),\quad
 \frac{11}{40}(4),\quad\frac{13}{40}(2),\quad
 \frac38(2),\quad\frac{17}{40}(2).
 \label{eq:spectrum}
\end{equation}
Consequently $W_*$ is a full-rank valid process and its smallest eigenvalue is $1/40$.

Let $\Pi_-=[I-\sum_{k=x,y,z}\sigma_k\otimes\sigma_k]/4$ be the singlet projector on $A_OB_I$. Define the Hermitian causal witness $S$ by
\begin{equation}
 S=(I-\sigma_z^{A_I})\Pi_-
   +\frac14(I-\sigma_z^{A_I})\sigma_x^{B_I}\sigma_z^{B_O}.
 \label{eq:witness}
\end{equation}
Its replaced operators are
\begin{align}
 {}_{B_O}S&=(I-\sigma_z^{A_I})\Pi_-\succeq0,\notag\\
 {}_{A_O}S&=\frac14(I-\sigma_z^{A_I})
        (I+\sigma_x^{B_I}\sigma_z^{B_O})\succeq0.
 \label{eq:witness-replacements}
\end{align}
For an $A\prec B$ summand $W^{A\prec B}$, self-adjointness of the replacement map gives
$\Tr(SW^{A\prec B})=\Tr[({}_{B_O}S)W^{A\prec B}]\ge0$; the other order uses \({}_{A_O}S\). Thus $S$ is a CS witness. Pauli orthogonality gives exactly
\begin{equation}
 \Tr(SW_*)=-\frac1{10},\qquad \Tr(SI/4)=1.
 \label{eq:witness-exact-values}
\end{equation}

\subsection{Relation to the existing transpose construction}
The ordinary causal model follows from the construction of Feix, Ara\'{u}jo, and Brukner~\cite{Feix}. In the original Bob port bases,
\begin{align}
 W_*^{T_{B_IB_O}}&=A_T+B_T,\notag\\
 A_T&=\frac1{20}\left[I+\sigma_x^{A_O}\sigma_x^{B_I}
 -\sigma_y^{A_O}\sigma_y^{B_I}+\sigma_z^{A_O}\sigma_z^{B_I}\right],\notag\\
 B_T&=\frac14\left[\frac45I+\frac15\sigma_z^{A_I}
 +\frac3{10}\sigma_z^{A_I}\sigma_x^{B_I}\sigma_z^{B_O}\right].
 \label{eq:FAB-relation}
\end{align}
Here $T_{B_IB_O}$ transposes only Bob's original input and output factors; the subscript $T$ on $A_T,B_T$ labels the components of this transposed process. $A_T$ is a positive $A\prec B$ component with trace $4/5$ and nonzero eigenvalue $1/5$. $B_T$ is a positive $B\prec A$ component with trace $16/5$ and minimum eigenvalue $3/40$. Both satisfy all process linear constraints. Bob's full Choi transpose preserves positive local instruments in the unextended experiment, giving an ordinary causal model for $W_*$.

This argument does not establish EC. Transposing an entire extended Bob instrument also transfers a partial transpose to the shared input state, which need not stay positive; transposing only the original ports need not preserve instrument positivity. Section VI and footnote 6 of Ref.~\cite{Feix} explicitly identify the gap between this construction, their finite searches, and extensible causality. The proof below supplies the missing arbitrary-ancilla statement.

\subsection{Operational response protocols}
\label{sec:responses}
The laboratory responses discussed in the main text follow directly from the Born rule, with no ancillary state required. The axis label $k\in\{x,y,z\}$ selects the Pauli matrix $\sigma_k$, independently of the setting label $x$. Let $\rho_x^k=(I+x\sigma_k)/2$ be its eigenstate of eigenvalue $x=\pm1$. Alice discards her input and prepares $\rho_x^k$; Bob measures $\sigma_k$, records $b=\pm1$, and returns a maximally mixed qubit. In the full-transpose Choi convention the instruments are
\begin{equation}
 M_x=I_{A_I}\otimes(\rho_x^k)^T,\qquad
 N_b=\frac{I+b\sigma_k}{2}\otimes\frac{I_{B_O}}2.
 \label{eq:forward-instruments}
\end{equation}
Here $M_x$ represents Alice's single-outcome preparation channel at setting $x$ and $N_b$ Bob's measurement event with outcome $b$. The resulting distribution is
\begin{equation}
 p(b|x,k)=\frac{1+\kappa_k x b/5}{2},\qquad
 \kappa_x=\kappa_z=1,\quad\kappa_y=-1.
 \label{eq:forward-response}
\end{equation}
Here $\kappa_k$ is the axis-dependent sign introduced by transposing the prepared state. The main-text protocol uses $k=x$. The $\sigma_y$ sign is checked separately to detect a Choi convention error.

For the reverse response, Bob measures $\sigma_x^{B_I}$ with result $r=\pm1$ and prepares a $\sigma_z^{B_O}$ eigenstate of eigenvalue $yr$, with setting $y=\pm1$. Alice measures $\sigma_z^{A_I}$ with result $s=\pm1$ and returns a maximally mixed qubit. Thus
\begin{equation}
 M_s=\frac{I+s\sigma_z}{2}\otimes\frac I2,\qquad
 N_{r|y}=\frac{I+r\sigma_x}{2}\otimes\frac{I+yr\sigma_z}{2}.
 \label{eq:reverse-instruments}
\end{equation}
Here $M_s$ represents Alice's event with measured result $s$, and $N_{r|y}$ Bob's measurement-and-preparation event with result $r$ and setting $y$; the tensor factors follow the laboratory input--output order. The joint distribution is $p(s,r|y)=[1+s/5+3sy/10]/4$. Summing over $r$ gives
\begin{equation}
 \langle b\rangle_x=\frac{x}{5}\quad(k=x),\qquad
 \langle s\rangle_y=\frac15+\frac{3y}{10}.
 \label{eq:directional-response}
\end{equation}
Here angle brackets denote the average recorded outcome conditional on the subscripted setting. Every displayed event is CP and each complete instrument is TP. These two separate experiments illustrate the operator's directional dependences; they do not supply its universal EC proof or a circuit realization. The independent verifier \texttt{repro/independent\_operational\_response\_audit.py} forms the standard Choi matrices from the physical maps, applies the full transpose, and checks all probabilities and a deliberately incorrect $\sigma_y$ transpose.

\section{Why nonseparability excludes every finite purification}
The pure bipartite process structure theorem~\cite{Yokojima,Cyclic} applies to a pure process whose unitary maps $P A_O B_O$ to $A_I B_I F$. It gives orthogonal decompositions of its past and future spaces into an $A\prec B$ sector and a $B\prec A$ sector. The unitary restricted to each sector is an ordered pure comb. This includes degenerate or zero sectors.

Insert any normalized pure past state $|p\rangle=|p_{A\prec B}\rangle+|p_{B\prec A}\rangle$, where the subscripts identify its components in the two ordered past sectors. Although these components may be coherent, their images lie in orthogonal future sectors. Tracing $F$ removes their cross terms. The reduced process is therefore
\begin{equation}
 W=W^{A\prec B}+W^{B\prec A},\qquad
 W^{A\prec B},W^{B\prec A}\succeq0,
 \label{eq:pure-to-CS}
\end{equation}
where $W^{A\prec B},W^{B\prec A}$ are the subnormalized ordered process components, with weights $\|p_{A\prec B}\|^2$ and $\|p_{B\prec A}\|^2$; these are the squared Hilbert-space norms of the past components. No rank condition or bound on the added environment enters this argument.

Conversely, a finite CS process has ordered-comb Stinespring dilations of its two normalized branches. Purify each branch's source and retain every channel environment. Padding the branch environments permits a controlled direct sum of the two ordered pure dilations. Prepare a coherent pure past state with amplitudes equal to the square roots of the branch weights and retain an orthogonal order flag in the future. Discarding that future yields the original CS process. Therefore, for the present scope of trivial \emph{original} past and future,
\begin{equation}
 \mathrm{Purifiable}=\CS.
 \label{eq:purifiable-CS}
\end{equation}
Here the equality concerns sets of processes within the stated scope. The witness~\eqref{eq:witness-exact-values} excludes $W_*$ from both sets. It remains to prove EC.

The same exclusion rules out a standard deterministic quantum circuit with quantum control of causal order (QC-QC). Every such circuit admits a unitary extension by retaining its internal environments~\cite{QCQC}. Preparing its fixed global input and discarding its future would constitute the purification excluded above. This statement concerns the standard circuit class in which each external operation is used once; postselected simulations and different implementation frameworks are outside its scope.

\section{Refinement and the complex Born reduction}
\subsection{Refining the two classical flags}
$W_*$ commutes with $\sigma_z^{A_I}$ and $\sigma_z^{B_O}$. Dephasing these ports in their $\sigma_z$ eigenbases leaves every Born probability unchanged. Refine Alice's report by its input flag $s=\pm1$ and Bob's report by its output flag $j=\pm1$. Sandwiching a Choi matrix with a flag projector preserves CP; summing refined reports preserves channel normalization. A causal model for the refined reports remains causal after their flags are discarded. We may hence prove EC for the refined experiment.

For a pure shared state choose Schmidt bases and write
\begin{equation}
 |\psi\rangle=\sum_i\sqrt{\rho_i}|i\rangle_{A'}|i\rangle_{B'},
 \qquad\rho=\sum_i\rho_i|i\rangle\langle i|,\qquad D=\sqrt\rho.
 \label{eq:schmidt}
\end{equation}
Here $|\psi\rangle$ is the shared pure state, $\rho_i$ its nonnegative Schmidt probabilities summing to one, $\rho$ the reduced density operator on Bob's paired Schmidt support, and $D$ its positive square root. The indices $i,j$ below label the paired ancillary basis vectors; $o,p\in\{0,1\}$ label Alice's output-qubit basis. Unused coordinates can be padded with zeros. An accepted subset of Alice's reports in input sector $s$ has Choi matrix $M_s$ on $A'A_O$. Its conditional physical CP output on $A_OB'$ has entries
\begin{equation}
 (\omega_s)_{o i,p j}=D_iD_j(M_s^T)_{i o,j p},\qquad
 \tau_s=\Tr_{A_O}\omega_s\preceq\rho,\qquad\omega_s\succeq0.
 \label{eq:sigma-choi}
\end{equation}
Here $\omega_s$ is the subnormalized joint output and $\tau_s$ its Bob marginal; $D_i=\sqrt{\rho_i}$ is the corresponding diagonal entry of $D$. The transpose in $M_s^T$ is the full transpose from Eq.~\eqref{eq:choi-convention}, followed by reordering to $A_OB'$. In particular,
\begin{equation}
 \tau_s=D(\Tr_{A_O}M_s)^T D\preceq D^2=\rho.
 \label{eq:tau-bound}
\end{equation}
$\omega_s$ itself is the untransposed physical CP output and can be mixed, with any finite number of Kraus branches. Write $\tau_{a,s|x}$ for the marginal associated with report $a$, sector $s$, and setting $x$. For the complete set of reports in either input sector, $\sum_a\tau_{a,s|x}=\rho$.

Bob's accepted subset in output sector $j$ gives an effect $E_j$ on $B_IB'$. Channel normalization implies
\begin{equation}
 E_+,E_-\succeq0,\qquad E=E_++E_-\preceq I_{B_IB'}.
 \label{eq:bob-effects}
\end{equation}
Here $E$ is the total effect of the accepted Bob subset, and $E_+,E_-$ its two output-flag components. For all refined reports together their effects sum to $I_{B_IB'}$.

\subsection{The effective Bob operator}
Under the identification of the $A_O$ and $B_I$ qubit bases, put $\omega_{s,k}=\Tr_{A_O}[(\sigma_k\otimes I)\omega_s]$ for $k=x,y,z$, where $\omega_{s,k}$ is the Pauli contraction of the joint output $\omega_s$. The process contraction contains $\sigma_k^T$, as follows directly from Eq.~\eqref{eq:sigma-choi}:
\begin{align}
 K_{s,j}={}&\frac{1+s\delta}{4}I\otimes\tau_s
     +\frac\beta4\sum_{k=x,y,z}\sigma_k^T\otimes\omega_{s,k}
     +\frac{s j\alpha}{4}\sigma_x\otimes\tau_s.\label{eq:kernel-pauli}
\end{align}
Here $K_{s,j}$ is the effective operator on Bob's input for Alice sector $s$ and Bob flag $j$. The Pauli sum equals $2\omega_s^{T_{A_O}}-I\otimes\tau_s$, where $T_{A_O}$ denotes partial transpose on Alice's output factor only. For a qubit,
\begin{equation}
 \omega_s^{T_{A_O}}=I\otimes\tau_s
                 -(\sigma_y\otimes I)\omega_s(\sigma_y\otimes I).
 \label{eq:qubit-transpose}
\end{equation}
Combining these identities gives
\begin{align}
 K_{s,j}={}&a_s I\otimes\tau_s
       -\frac\beta2(\sigma_y\otimes I)\omega_s(\sigma_y\otimes I)
       +\frac{s j\alpha}{4}\sigma_x\otimes\tau_s,\notag\\
 a_+={}&\frac7{20},\qquad a_- =\frac14.
 \label{eq:kernel}
\end{align}
Here $a_s$ is the flag-dependent scalar coefficient of $I\otimes\tau_s$. The probability of the selected pair of subsets is $\sum_{s,j}\Tr(E_jK_{s,j})$. This derivation fixes the $\sigma_y$ sign and the Choi conjugations for complex instruments. An extra transpose of $\omega_s$, of the shared state, or of a Bob block changes the contraction.

\section{The transportation criterion for all event subsets}
Choose the same constants for every experiment:
\begin{equation}
 q=\frac{49}{80},\quad\theta=\frac3{25},\quad
 w_s=\frac{q+s\theta}{2},\quad
 w_+=\frac{293}{800},\quad w_-=\frac{197}{800},\quad
 c=\frac{1-q}{2}=\frac{31}{160}.
 \label{eq:margins}
\end{equation}
Here $q$ is the order weight, $\theta$ adjusts Alice's relative flag weights $w_s$, and $c$ is the common coefficient of Bob's reference contribution. For refined reports $a$ and $b$, with the port flags included in the reports, prescribe
\begin{equation}
 g_a(x)=w_s\Tr\tau_{a,s|x},\qquad
 h_b(y)=c\Tr[(I_{B_I}\otimes\rho)E_{b|y}].
 \label{eq:prescribed-margins}
\end{equation}
Here $g_a(x),h_b(y)$ are the prescribed subnormalized first-party marginals in the two ordered branches and $E_{b|y}$ is Bob's refined event effect. These marginals are nonnegative and depend only on their own party's setting. Their total masses are $\sum_a g_a=q$ and $\sum_b h_b=2c=1-q$.

\begin{lemma}[Capacitated transportation]
For a normalized nonnegative table $p_{ab}$ and nonnegative $g_a,h_b$ of masses $q,1-q$, respectively, there are nonnegative tables $u,v$ such that $p=u+v$, $\sum_bu_{ab}=g_a$, and $\sum_av_{ab}=h_b$ if and only if
\begin{equation}
 p(\mathcal A,\mathcal B)\le g(\mathcal A)+h(\mathcal B)
 \quad\text{for every pair of outcome subsets }\mathcal A,\mathcal B.
 \label{eq:hall}
\end{equation}
\end{lemma}
Here $p_{ab}$ is the joint table at a fixed setting pair. $p(\mathcal A,\mathcal B)$ sums its entries over $a\in\mathcal A,b\in\mathcal B$, while $g(\mathcal A),h(\mathcal B)$ sum the prescribed marginals over their respective subsets. The tables $u,v$ are the subnormalized ordered contributions.
\begin{proof}
Necessity follows from $u(\mathcal A,\mathcal B)\le g(\mathcal A)$ and $v(\mathcal A,\mathcal B)\le h(\mathcal B)$. For sufficiency, including full and complementary subsets in Eq.~\eqref{eq:hall} gives $p_A(a)\ge g_a$ and $p_B(b)\ge h_b$, where $p_A,p_B$ are the row and column marginals of $p$. Let $k_b=p_B(b)-h_b\ge0$ be the residual column demand, of total mass $q$. A network with source--row capacities $g_a$, row--column capacities $p_{ab}$, and column--sink capacities $k_b$ has a flow of mass $q$ exactly when each cut has capacity at least $q$. If its source-side rows and columns are $\mathcal S,\mathcal T$, that condition is
\begin{align}
 &g(\mathcal S^c)+p(\mathcal S,\mathcal T^c)+k(\mathcal T)-q\notag\\
 &\hspace{15mm}=g(\mathcal S^c)+h(\mathcal T^c)
                   -p(\mathcal S^c,\mathcal T^c)\ge0.
 \label{eq:flow-cuts}
\end{align}
Here the superscript $c$ denotes the complement within the corresponding outcome set and $k(\mathcal T)$ sums the residual demands over $\mathcal T$. These are precisely the assumed inequalities. The flow entries define $0\le u_{ab}\le p_{ab}$, with row sums $g_a$ and column sums $k_b$; setting $v=p-u$ proves the claim.
\end{proof}

Apply the lemma independently to every setting pair $(x,y)$. The same prescribed $g_a(x)$ and $h_b(y)$ are used throughout. Then $u/q$ is normalized and has Alice marginal $g_a(x)/q$, independent of $y$. Likewise $v/(1-q)$ has Bob marginal $h_b(y)/(1-q)$, independent of $x$. The setting-independent mixture weight is the fixed $q$. Thus proving every inequality~\eqref{eq:hall} proves causality of the whole experiment, regardless of the numbers of settings or outcomes.

Using the Born reduction, the left side minus the right side of an arbitrary such inequality is
\begin{equation}
 H=\sum_{s,j}\Tr(E_jK_{s,j})
       -\sum_s w_s\Tr\tau_s-c\Tr[(I\otimes\rho)E].
 \label{eq:Hall-score}
\end{equation}
Here $H$ is the scalar Hall excess for the selected event subsets. Sections~\ref{sec:lifting}--\ref{sec:universal} prove $H\le0$ for all allowed events and states, including empty, full, and complementary subsets.

\section{Lifting general CP events and Bob effects}
\label{sec:lifting}
\subsection{Retaining arbitrary Kraus environments}
Write Alice's CP event in sector $s$ with Kraus matrices $K_{s,\mu}:\mathcal H_A\to\mathbb C^2$, where $\mathcal H_A$ is her finite ancillary input space, $\mathbb C^2$ her qubit output space, and $\mu$ the Kraus-branch label. They satisfy $\sum_\mu K_{s,\mu}^\dagger K_{s,\mu}\preceq I_A$, with $I_A$ the identity on $\mathcal H_A$. Retain their environment and set
\begin{equation}
 L_{s,o}=\sum_\mu|\mu\rangle\langle o|K_{s,\mu},\qquad
 R_s=I_A-\sum_{o=0}^1L_{s,o}^\dagger L_{s,o}\succeq0.
 \label{eq:Alice-contraction}
\end{equation}
Here $|\mu\rangle$ is an orthonormal basis of the retained environment $\mathcal H_E$, $\langle o|$ selects output basis state $o$, and $R_s$ is the positive slack left by the accepted event. The rectangular amplitudes $L_{s,o}:\mathcal H_A\to\mathcal H_E$ reproduce $\omega_s$ after the environment is traced. Zero Kraus matrices pad the two environments to a common finite space. To interpret a square-matrix polynomial literally, embed $\mathcal H_A$ and $\mathcal H_E$ in a common larger Alice space and pad each $L_{s,o}$ with zero rows and columns. The shared state is embedded in the original input subspace. This preserves all $L_{s,o}^\dagger L_{s,p}$ matrix elements and leaves $I-\sum_oL_{s,o}^\dagger L_{s,o}\succeq0$ on the whole enlarged space. Therefore arbitrary CP events, including mixed outputs and multiple Kraus branches, are covered by the contraction constraints.

\subsection{A dilation preserving the input-qubit operators}
Set $E_0=I-E_+-E_-$ for the rejected effect. This three-outcome positive operator-valued measure (POVM) has a finite Naimark dilation. It can be chosen with the fixed input embedding
\begin{equation}
 J=I_{B_I}\otimes I_{B'}\otimes|0\rangle,\qquad
 J^\dagger P_jJ=E_j,\qquad P_jP_\ell=\delta_{j\ell}P_j,
 \label{eq:Naimark}
\end{equation}
where $J$ appends a fixed ancillary basis state $|0\rangle$, $P_j$ are the dilated event projectors, and $j,\ell\in\{+,-,0\}$ label the three outcomes. The Kronecker delta $\delta_{j\ell}$ is one for equal labels and zero otherwise; $\sum_jP_j=I$. To construct the projectors, let $|j\rangle$ be the three orthonormal ancillary outcome vectors. The isometry $V_N=\sum_j\sqrt{E_j}\otimes|j\rangle$ extends to a unitary $U_N$ with $U_NJ=V_N$; take $P_j=U_N^\dagger(I\otimes|j\rangle\langle j|)U_N$.

The added input space belongs to Bob's ancilla; $B_I$ remains dimension two. The embedding $J$ intertwines every $B_I$ Pauli matrix. Thus Eq.~\eqref{eq:Naimark} preserves the $I$, $\sigma_x$, and $\sigma_y$ terms of Eq.~\eqref{eq:kernel}, as well as the cost $\Tr[(I\otimes\rho)E_j]$. It preserves $H$ exactly. We can hence replace Bob's effects by $P_j$ and the shared state by its fixed local embedding.

\subsection{The Hall operator}
For the dilated operators define the unnormalized qubit partial traces and blocks
\begin{equation}
 T=\Tr_{B_I}(P_++P_-),\qquad
 V=\Tr_{B_I}[\sigma_x(P_+-P_-)],\qquad
 F=\sigma_y(P_++P_-)\sigma_y.
 \label{eq:TVF}
\end{equation}
Here $T$ is the total accepted effect traced over the input qubit, $V$ its $\sigma_x$-weighted outcome difference, and $F$ the $\sigma_y$-conjugated accepted effect on the qubit and enlarged ancilla. In this subsection $F$ is an operator, distinct from the global future space in Section S4. The Pauli matrices in Eq.~\eqref{eq:TVF} act on $B_I$, with identities on Bob's enlarged ancilla. The $(o,p)$ qubit block is denoted by $F_{op}$. All Alice operators commute with all Bob operators. Define the Hall operator by
\begin{align}
 \mathcal H={}&\sum_{s,o}L_{s,o}^\dagger L_{s,o}
          (a_sT+s\alpha V/4-w_sI)\notag\\
 &-\frac\beta2\sum_{s,o,p}L_{s,o}^\dagger L_{s,p}F_{op}-cT
 \label{eq:operator-H}
\end{align}
where $\mathcal H$ acts on the two enlarged auxiliary spaces and its shared-state expectation is exactly the scalar $H$. In particular, the indices in its second line are $F_{op}$: expanding $\Tr(F\omega_s)$ gives $L_{s,o}^\dagger L_{s,p}F_{op}$. No transpose of this block occurs.

\section{A complete rational SOS identity}
\label{sec:sos}
All the following matrices have operator-valued entries on the two auxiliary spaces, and their sizes are independent of those spaces' dimensions. For a real symmetric scalar matrix $N$ and a compatible operator matrix $Q$, define
\begin{equation}
 \langle N,Q\rangle=\sum_{i,j}N_{ij}Q_{ij}.
 \label{eq:weighted-contraction}
\end{equation}
Here the indices $i,j$ run over the common matrix size; the angle brackets here denote a weighted matrix contraction, rather than the outcome average of Section S3.2. For $N,Q\succeq0$ this contraction is positive. It is the scalar-index partial trace of $N^{1/2}QN^{1/2}$.

\subsection{Positive matrices and their bases}
Write $C=P_0=I-P_+-P_-$ for the rejected projector, with qubit row and column indices $r,c\in\{0,1\}$ and blocks $C_{rc}$. Use row-major ordering throughout and define the operator tuples
\begin{align}
 d&=(C_{00},C_{01},C_{10},C_{11}),\notag\\
 n&=(C_{00}-C_{11},C_{01},C_{10}),\notag\\
 e&=((P_+)_{00},(P_+)_{01},(P_+)_{10},(P_+)_{11},
       (P_-)_{00},(P_-)_{01},(P_-)_{10},(P_-)_{11}),\notag\\
 v_s&=(L_{s,0}d,L_{s,1}d).
 \label{eq:SOS-bases}
\end{align}
Products of an amplitude with a tuple act on each entry, and parentheses concatenate the resulting tuples. For an operator tuple $z=(z_i)$ define the Gram and slack-localized Gram matrices by
\begin{equation}
 G(z)_{ij}=z_i^\dagger z_j,\qquad
 \mathcal L_{R_s}(z)_{ij}=z_i^\dagger R_s z_j.
 \label{eq:SOS-gram}
\end{equation}
Here $G(z)$ is the Gram matrix and $\mathcal L_{R_s}(z)$ inserts the positive Alice slack between tuple entries. Both are positive. The remaining two types are
\begin{equation}
 (\mathcal M_s)_{(o,r),(p,c)}=L_{s,o}^\dagger C_{rc}L_{s,p},\qquad
 (D_{s,j})_{rc}=R_s(P_j)_{rc}.
 \label{eq:SOS-localizers}
\end{equation}
Here $\mathcal M_s$ is the accepted-amplitude localizer of the rejected projector and $D_{s,j}$ the slack--projector product matrix. Positivity of $\mathcal M_s$ follows by applying $C\succeq0$ to $\sum_oL_{s,o}z_{o,r}$ for arbitrary test vectors $z_{o,r}$. Positivity of $D_{s,j}$ follows because $R_s$ and $P_j$ act on different laboratories. Compress $\mathcal M_s$ to
\begin{equation}
 \widetilde{\mathcal M}_s=B^T\mathcal M_sB,\qquad
 B=\begin{pmatrix}1&0&0\\0&1&0\\0&1&0\\0&0&1\end{pmatrix},
 \label{eq:SOS-triplet}
\end{equation}
where $B$ is the displayed real $4\times3$ compression matrix, $\widetilde{\mathcal M}_s$ the compressed localizer, and the four rows correspond to $(o,r)=(0,0),(0,1),(1,0),(1,1)$.

The eight-entry bases split under a reflection. For $v_s$ choose representative and partner indices $(0,1,2,3)$ and $(7,6,5,4)$; for $e$ choose $(0,1,4,5)$ and $(3,2,7,6)$. Indices start at zero. Define four-entry tuples
\begin{equation}
 u_s^{\mathrm e/\mathrm o}=(v_s)_{\mathrm{rep}}\mathbin{\pm}(v_s)_{\mathrm{partner}},\qquad
 z^{\mathrm e/\mathrm o}=e_{\mathrm{rep}}\mathbin{\pm}e_{\mathrm{partner}}.
 \label{eq:SOS-parity}
\end{equation}
Here ``rep'' and ``partner'' extract entries in the ordered index lists just specified; superscripts $\mathrm e,\mathrm o$ denote even and odd reflection parity. The even tuple uses $+$ and the odd tuple uses $-$. There is no square-root normalization factor.

\subsection{The fifteen scalar contractions}
Let $\kappa=4\,000\,000$ be the common weight denominator and $k_+=1,k_-=8$ the starting matrix indices for the two flag sectors. The fifteen integer coefficient matrices $N_0,\ldots,N_{14}$ listed below define the unscaled positive certificate operator $\mathcal Q$ by
\begin{align}
 \mathcal Q={}&\langle N_0,G(n)\rangle\notag\\
 &+\sum_{s=+,-}\Bigl[
    \langle N_{k_s},G(u_s^{\mathrm e})\rangle
   +\langle N_{k_s+1},G(u_s^{\mathrm o})\rangle\notag\\
 &\qquad+\langle N_{k_s+2},\mathcal L_{R_s}(z^{\mathrm e})\rangle
   +\langle N_{k_s+3},\mathcal L_{R_s}(z^{\mathrm o})\rangle\notag\\
 &\qquad+\langle N_{k_s+4},\widetilde{\mathcal M}_s\rangle
   +\langle N_{k_s+5},D_{s,+}\rangle
   +\langle N_{k_s+6},D_{s,-}\rangle\Bigr].
 \label{eq:compact-sos}
\end{align}
Their actual SOS weights are $N_k/\kappa$. The explicit matrices are:
% Generated by repro/export_hall_sos_tables.py; all matrices are exact.
% Every displayed N_k corresponds to the weight N_k/4000000.
\begingroup
\small
\allowdisplaybreaks
% scalar_Bob_traceless_complement_Gram; basis: identity
\begin{equation*}
N_{0}=\begin{pmatrix}
3368 & 0 & 0 \\
0 & 3370 & 0 \\
0 & 0 & 3370
\end{pmatrix}
\end{equation*}
% flag_0_forward_complement_Gram_even; basis: u_i=v_rep_i+v_partner_i
\begin{equation*}
N_{1}=\begin{pmatrix}
376525 & 0 & 0 & 40190 \\
0 & 164200 & 0 & 178330 \\
0 & 0 & 336335 & 0 \\
40190 & 178330 & 0 & 204390
\end{pmatrix}
\end{equation*}
% flag_0_forward_complement_Gram_odd; basis: u_i=v_rep_i-v_partner_i
\begin{equation*}
N_{2}=\begin{pmatrix}
376525 & 0 & 0 & 40190 \\
0 & 164200 & 0 & -178330 \\
0 & 0 & 336335 & 0 \\
40190 & -178330 & 0 & 204390
\end{pmatrix}
\end{equation*}
% flag_0_Alice_slack_Bob_Gram_even; basis: u_i=v_rep_i+v_partner_i
\begin{equation*}
N_{3}=\begin{pmatrix}
191947 & 109881 & 80380 & 0 \\
109881 & 111567 & 0 & 0 \\
80380 & 0 & 191947 & -109881 \\
0 & 0 & -109881 & 111567
\end{pmatrix}
\end{equation*}
% flag_0_Alice_slack_Bob_Gram_odd; basis: u_i=v_rep_i-v_partner_i
\begin{equation*}
N_{4}=\begin{pmatrix}
111567 & 109881 & 0 & 0 \\
109881 & 111567 & 0 & 0 \\
0 & 0 & 111567 & -109881 \\
0 & 0 & -109881 & 111567
\end{pmatrix}
\end{equation*}
% flag_0_Bob_complement_localizer; basis: identity
\begin{equation*}
N_{5}=\begin{pmatrix}
99070 & 0 & 0 \\
0 & 43340 & 0 \\
0 & 0 & 99070
\end{pmatrix}
\end{equation*}
% flag_0_Bob_label_0_slack_product; basis: identity
\begin{equation*}
N_{6}=\begin{pmatrix}
83606 & 80238 \\
80238 & 83606
\end{pmatrix}
\end{equation*}
% flag_0_Bob_label_1_slack_product; basis: identity
\begin{equation*}
N_{7}=\begin{pmatrix}
83606 & -80238 \\
-80238 & 83606
\end{pmatrix}
\end{equation*}
% flag_1_forward_complement_Gram_even; basis: u_i=v_rep_i+v_partner_i
\begin{equation*}
N_{8}=\begin{pmatrix}
383070 & 0 & 0 & -38506 \\
0 & 223259 & 0 & 198314 \\
0 & 0 & 423259 & 0 \\
-38506 & 198314 & 0 & 183070
\end{pmatrix}
\end{equation*}
% flag_1_forward_complement_Gram_odd; basis: u_i=v_rep_i-v_partner_i
\begin{equation*}
N_{9}=\begin{pmatrix}
383070 & 0 & 0 & -38506 \\
0 & 223259 & 0 & -198314 \\
0 & 0 & 423259 & 0 \\
-38506 & -198314 & 0 & 183070
\end{pmatrix}
\end{equation*}
% flag_1_Alice_slack_Bob_Gram_even; basis: u_i=v_rep_i+v_partner_i
\begin{equation*}
N_{10}=\begin{pmatrix}
110330 & -108517 & 113 & 0 \\
-108517 & 189025 & 0 & 78808 \\
113 & 0 & 110330 & 108517 \\
0 & 78808 & 108517 & 189025
\end{pmatrix}
\end{equation*}
% flag_1_Alice_slack_Bob_Gram_odd; basis: u_i=v_rep_i-v_partner_i
\begin{equation*}
N_{11}=\begin{pmatrix}
187342 & -108517 & 77125 & 0 \\
-108517 & 189025 & 0 & 78808 \\
77125 & 0 & 187342 & 108517 \\
0 & 78808 & 108517 & 189025
\end{pmatrix}
\end{equation*}
% flag_1_Bob_complement_localizer; basis: identity
\begin{equation*}
N_{12}=\begin{pmatrix}
3372 & 0 & 0 \\
0 & 3372 & 0 \\
0 & 0 & 3372
\end{pmatrix}
\end{equation*}
% flag_1_Bob_label_0_slack_product; basis: identity
\begin{equation*}
N_{13}=\begin{pmatrix}
86840 & -82966 \\
-82966 & 86840
\end{pmatrix}
\end{equation*}
% flag_1_Bob_label_1_slack_product; basis: identity
\begin{equation*}
N_{14}=\begin{pmatrix}
86840 & 82966 \\
82966 & 86840
\end{pmatrix}
\end{equation*}
\endgroup


The matrix pairing and bases in Eqs.~\eqref{eq:SOS-bases}--\eqref{eq:compact-sos} specify each scalar coefficient. For example, $N_1,N_2$ weight Alice's $s=+$ forward Gram tuples, $N_3,N_4$ the same sector's slack-localized Bob tuples, $N_5$ its compressed complement localizer, and $N_6,N_7$ its two slack products. $N_8$ through $N_{14}$ have the corresponding roles for $s=-$.

\subsection{All projection-ideal multipliers}
The only equalities used in the identity are Bob's projector block relations. For $j,\ell\in\{+,-\}$ and $r,c\in\{0,1\}$, denote the relation residuals by
\begin{equation}
 \mathfrak r^{j\ell}_{rc}=
   \sum_{t=0}^1(P_j)_{rt}(P_\ell)_{tc}
            -\delta_{j\ell}(P_j)_{rc}=0.
 \label{eq:ideal-relations}
\end{equation}
There are sixteen such block polynomials $\mathfrak r^{j\ell}_{rc}$, all zero for the physical projectors. Their scalar and Alice-quadratic multiples give all 144 coefficients of the certificate. Define the ideal contribution $\mathcal I$ and its operator multipliers $B^{j\ell}_{rc}$ by
\begin{align}
 \mathcal I&=\sum_{j,\ell,r,c}B^{j\ell}_{rc}\,\mathfrak r^{j\ell}_{rc},\notag\\
 B^{j\ell}_{rc}&=b^{j\ell}_{rc}I+
 \sum_{s,o,p}(C_s^{j\ell})_{(o,r),(p,c)}
                 L_{s,o}^\dagger L_{s,p}.
 \label{eq:ideal-coefficients}
\end{align}
Here $b^{j\ell}_{rc}$ are scalar coefficients and $C_s^{j\ell}$ their $4\times4$ coefficient matrices for Alice's quadratic amplitudes; all entries are integers. Write $b^{j\ell}$ for the $2\times2$ matrix with entries $b^{j\ell}_{rc}$, and introduce its integer scalar coefficients $a_0,b_0,c_0$:
\begin{equation}
 b^{++}=a_0I+b_0\sigma_x,\qquad b^{--}=a_0I-b_0\sigma_x,\qquad
 b^{+-}=b^{-+}=c_0I.
 \label{eq:ideal-bare}
\end{equation}
For each $s$, rows and columns of $C_s^{j\ell}$ are ordered $(o,r)=(0,0),(0,1),(1,0),(1,1)$.

Let $Z_4=\operatorname{diag}(1,-1,-1,1)$ be the diagonal reflection matrix. The numerical coefficient matrices $U_s,V_s$ determine the four label combinations through
\begin{equation}
 C_s^{++}=U_s,\quad C_s^{+-}=V_s,\quad
 C_s^{--}=Z_4U_sZ_4,\quad C_s^{-+}=V_s^T=Z_4V_sZ_4.
 \label{eq:ideal-symmetry}
\end{equation}
The constants and four integer matrices are:
% Generated by repro/export_hall_sos_tables.py; compact projection-ideal coefficients.
% All coefficient matrices share denominator 4000000.
\begin{equation*}
a_0=-604554,\qquad b_0=-2728,\qquad c_0=-160986.
\end{equation*}
\begin{equation*}
U_{+}=\begin{pmatrix}
-449536 & 219762 & 0 & 0 \\
219762 & -105266 & -356660 & 0 \\
0 & -356660 & -105266 & 219762 \\
0 & 0 & 219762 & -449536
\end{pmatrix}
\end{equation*}
\begin{equation*}
V_{+}=\begin{pmatrix}
-672670 & 0 & 0 & 0 \\
0 & -328400 & -356660 & 0 \\
0 & -356660 & -328400 & 0 \\
0 & 0 & 0 & -672670
\end{pmatrix}
\end{equation*}
\begin{equation*}
U_{-}=\begin{pmatrix}
-468468 & -217034 & 0 & 0 \\
-217034 & -68468 & -396628 & 0 \\
0 & -396628 & -68468 & -217034 \\
0 & 0 & -217034 & -468468
\end{pmatrix}
\end{equation*}
\begin{equation*}
V_{-}=\begin{pmatrix}
-688902 & 0 & 0 & 0 \\
0 & -288902 & -396628 & 0 \\
0 & -396628 & -288902 & 0 \\
0 & 0 & 0 & -688902
\end{pmatrix}
\end{equation*}


Equations~\eqref{eq:ideal-bare}--\eqref{eq:ideal-symmetry} thus print every ideal multiplier without leaving a free coefficient. Dividing $B^{j\ell}_{rc}$ by $\kappa$ gives the multipliers for the weights $N_k/\kappa$.

\subsection{The exact identity and positivity}
With the rational parameters~\eqref{eq:margins}, direct expansion gives the polynomial identity
\begin{equation}
 \boxed{-\kappa\mathcal H=\mathcal Q+\mathcal I.}
 \label{eq:SOS-identity}
\end{equation}
Here $\mathcal I$ vanishes under the physical projector relations, while $\mathcal Q$ is positive. The identity holds in free Alice and Bob words with interparty commutation and the usual adjoint rule; within each party the multiplication order is retained. Every $N_k$ is real symmetric and strictly positive definite. Table~\ref{tab:hall-sos-leading-minors} gives their exact leading principal determinants $\Delta_r(N_k)$, with $r$ the principal-block size. In an LDL factorization, $L$ is unit lower triangular and $D$ diagonal; the positive pivots $d_r$ are successive determinant ratios. Consequently $\mathcal Q\succeq0$ and $\mathcal H\preceq0$.
% Generated by repro/export_hall_sos_tables.py using exact integer determinants.
\begin{table}[t]
\centering
\scriptsize
\setlength{\tabcolsep}{3pt}
\begin{tabular}{r r r r r}
\hline
$k$ & $\Delta_1(N_k)$ & $\Delta_2(N_k)$ & $\Delta_3(N_k)$ & $\Delta_4(N_k)$ \\
\hline
0 & 3368 & 11350160 & 38250039200 & $-$ \\
1 & 376525 & 61825405000 & 20794047590675000 & 133585367050938762500 \\
2 & 376525 & 61825405000 & 20794047590675000 & 133585367050938762500 \\
3 & 191947 & 9341116788 & 1072171160231436 & 6835824856695625344 \\
4 & 111567 & 373361328 & 41654803280976 & 139398681245923584 \\
5 & 99070 & 4293693800 & 425376244766000 & $-$ \\
6 & 83606 & 551826592 & $-$ & $-$ \\
7 & 83606 & 551826592 & $-$ & $-$ \\
8 & 383070 & 85523825130 & 36198728700698670 & 110162193384172888504 \\
9 & 383070 & 85523825130 & 36198728700698670 & 110162193384172888504 \\
10 & 110330 & 9079188961 & 1001704504406905 & 7039978207917249424 \\
11 & 187342 & 23636382261 & 3303716215774637 & 308260026331361563600 \\
12 & 3372 & 11370384 & 38340934848 & $-$ \\
13 & 86840 & 657828444 & $-$ & $-$ \\
14 & 86840 & 657828444 & $-$ & $-$ \\
\hline
\end{tabular}
\caption{Exact leading principal determinants of the fifteen integer SOS matrices $N_k$. The actual weights are $N_k/4000000$, so each $r$th leading determinant is $\Delta_r(N_k)/4000000^r$. All displayed integers are strictly positive; Sylvester's criterion proves positive definiteness. A dash denotes a principal block larger than the matrix. The LDL pivots are $d_r=\Delta_r/\Delta_{r-1}$ with $\Delta_0=1$.}
\label{tab:hall-sos-leading-minors}
\end{table}


The eleven-block certificate before reflection splitting is also supplied. Its integer coefficient system has size $656\times318$, rank $271$, and $47$ free coordinates. Rational recovery satisfies all coefficient equations exactly, and exact LDL factorizations certify all eleven weights. Averaging under joint $A_O/B_I$ Pauli $\sigma_x$, then under joint Pauli $\sigma_z$ with Bob's $+/-$ reports interchanged, preserves the Hall polynomial and positivity. Splitting the four eight-dimensional blocks yields precisely the fifteen matrices printed here. This transformation is algebraic and uses no further numerical optimization.

No Alice common-range, orthogonality, or nilpotent relation is used in Eq.~\eqref{eq:SOS-identity}. Its only Alice assumption is $R_s\succeq0$. The identity is therefore valid for arbitrary finite auxiliary dimensions. Taking its expectation in the embedded shared state proves $H\le0$.

\section{Universal EC and an independent conic route}
\label{sec:universal}
The transportation lemma gives a useful sufficient criterion beyond the particular coefficients of $W_*$. For each shared state, prescribe nonnegative first-party marginals $g_a(x)$ and $h_b(y)$ of masses $q$ and $1-q$, where $q$ is independent of the settings. If every physical event-subset pair has Hall excess at most zero for all settings, the entire experiment is causal. An operator SOS valid in all physical representations is one way of proving these bounds simultaneously. The transportation criterion is general; the kernel, marginal choice, and rational weights supplied here certify the explicit example.

The arbitrary-dimensional inequality applies to the two Alice event subsets associated with any $\mathcal A$, and the two Bob event effects associated with any $\mathcal B$. Thus it proves every cut in Eq.~\eqref{eq:hall}. The transportation lemma proves a causal decomposition for every pure shared state and every finite instrument collection. For a mixed shared state, decompose it into pure states and combine their causal models; the weight $q$ in Eq.~\eqref{eq:margins} is common to all components. Equivalently, purify the state into an enlarged local input ancilla. Refinement and coarse graining then prove
\begin{equation}
 W_*\in\EC\setminus\CS,\qquad W_*\notin\mathrm{Purifiable}.
 \label{eq:final-theorem}
\end{equation}
Here $\EC\setminus\CS$ means the set of EC processes excluding CS processes, $\in$ denotes membership, and $\notin$ nonmembership.

A second exact certificate is retained for the original, larger canonical conic relaxation. Its exact dual equations and zero objective provide an independent route to the same Hall bound. For completeness, the additional physical representation used by that relaxation can be attained without restricting CP events. Complete each accepted event by rejected Kraus matrices, pad the accepted slots for both sectors to the same number $a$, and pad the total environment to dimension $m$. Let $Q_E$ project onto these $a$ accepted slots. The complete Stinespring isometry $\mathcal V_s:\mathcal H_A\to\mathbb C^2\otimes\mathbb C^m$ extends to a unitary $\mathcal U_s$ on a space of dimension $2m$, with a common input embedding $J_A$ and $\mathcal U_sJ_A=\mathcal V_s$. For any common isometry $F_E:\mathbb C^m\to\mathbb C^{2m}$ set
\begin{equation}
 D_{s,o}=F_EQ_E(\langle o|\otimes I_m)\mathcal U_s,\qquad
 Q=F_EQ_EF_E^\dagger.
 \label{eq:common-range-dilation}
\end{equation}
Here $D_{s,o}$ are dilated accepted-event amplitudes, $Q$ their common output-range projector, and $I_m$ the identity on the $m$-dimensional environment. These matrices are distinct from the Schmidt square root $D$ and the slack products $D_{s,j}$ above. Then $D_{s,o}D_{s,p}^\dagger=\delta_{op}Q$ and
$\sum_oD_{s,o}^\dagger D_{s,o}=\mathcal U_s^\dagger(I_2\otimes Q_E)\mathcal U_s$ is a projector. Tracing the retained output environment recovers every original accepted CP output, including its mixedness. The Bob marginal and all Born pairings are preserved.

One further doubling gives
\begin{equation}
 A_{s,o}=\begin{pmatrix}0&0\\D_{s,o}&0\end{pmatrix},\qquad
 J_2=\begin{pmatrix}J_A\\0\end{pmatrix}.
 \label{eq:nilpotent-dilation}
\end{equation}
Here $A_{s,o}$ are the amplitudes embedded in the doubled space and $J_2$ its input embedding. This preserves the common output range and enforces
$A_{s,o}A_{t,p}=0$ and $A_{s,o}^\dagger J_2=0$ for all sector labels $s,t$ and output-basis indices $o,p$. This gives the state-supported 45-word Gram basis of the canonical relaxation. These relations apply to this representation, not to arbitrary original event matrices. The small certificate of Section~\ref{sec:sos} does not need this construction.

The independent rational verifier checks $10\,960$ canonical dual equations, exact upper bound zero, and positivity of its supported cone blocks. It separately reconstructs the small SOS from the original Pauli contraction, finding $48$ nonzero target words and zero residual words. The complete finite-dimensional proof above supplies physical coverage and all transportation cuts. Numerical examples are convention and implementation checks rather than a substitute for this proof.

\section{Exact classification on the white-noise line}
Let
\begin{equation}
 W_\eta=(1-\eta)W_*+\eta I/4,\qquad 0\le\eta\le1.
 \label{eq:noise}
\end{equation}
Here $\eta$ is the probability of replacing the counterexample by the uniform process $I/4$, and the subscript labels the resulting noisy process. EC is convex, and $I/4$ is CS and hence EC. Therefore $W_\eta$ is EC everywhere on this line. The witness value is
\begin{equation}
 \Tr(SW_\eta)=\frac{11\eta-1}{10}.
 \label{eq:noise-witness}
\end{equation}
It excludes CS and purification for $0\le\eta<1/11$.

To prove exactness, set $v=1-\eta$ for the remaining visibility and write the two ordered components
\begin{align}
 W_\eta^{A\prec B}&=\frac14\left[\frac{3v}{5}I
          +\frac v5\sum_{k=x,y,z}\sigma_k^{A_O}\sigma_k^{B_I}\right],\notag\\
 W_\eta^{B\prec A}&=\frac14\left[\left(1-\frac{3v}{5}\right)I
          +\frac v5\sigma_z^{A_I}
          +\frac{3v}{10}\sigma_z^{A_I}\sigma_x^{B_I}\sigma_z^{B_O}\right].
 \label{eq:noise-ordered-decomposition}
\end{align}
Their sum is $W_\eta$, and each obeys its ordered-process linear constraints. The first is positive with eigenvalues $0$ and $v/5$ and weight $3v/5$. The second has weight $1-3v/5$ and minimum eigenvalue $(1-11v/10)/4$, which is nonnegative exactly when $\eta\ge1/11$. Hence
\begin{equation}
 W_\eta\in\EC\quad\text{for all }\eta,\qquad
 W_\eta\in\mathrm{Purifiable}\Longleftrightarrow
 W_\eta\in\CS\Longleftrightarrow\eta\ge\frac1{11}.
 \label{eq:noise-classification}
\end{equation}
Here the double arrows mean ``if and only if.'' The boundary $\eta=1/11$ is included in CS and purifiability. The line admits the following stability extension in the full valid-process space.

\section{A neighborhood stable in every valid perturbation direction}
\label{sec:robustness}
All norms and traces in this section use the process normalization $\Tr W=4$. Let $\mathcal V_0$ be the Hermitian, trace-zero subspace obeying the homogeneous constraints~\eqref{eq:process-subspace}; its real dimension is $87$. For a Hermitian perturbation $D$ in this section, the operator norm $\|D\|_\infty$ is its largest absolute eigenvalue; this $D$ is distinct from the Schmidt square root in Section S5. The trace norm $\|S\|_1$ is the sum of absolute eigenvalues of the witness.

\begin{lemma}[An explicit CS inner ball]
For every $D\in\mathcal V_0$,
\begin{equation}
 \|D\|_\infty\le\frac1{12}
 \quad\Longrightarrow\quad I/4+D\in\CS.
 \label{eq:CS-ball}
\end{equation}
\end{lemma}
\begin{proof}
Put $R_A={}_{A_O}$ and $R_B={}_{B_O}$. These commuting output-replacement maps preserve the process linear subspace. They are positive and unital, so they contract the operator norm on Hermitian inputs: apply each map to $-\|D\|_\infty I\preceq D\preceq\|D\|_\infty I$. Equation~\eqref{eq:process-subspace} gives $D=R_AD+R_BD-R_AR_BD$. Define
\begin{align}
 A(D)&=\frac I8+R_BD-\frac12R_AR_BD,\notag\\
 B(D)&=\frac I8+R_AD-\frac12R_AR_BD.
 \label{eq:CS-ball-split}
\end{align}
Here $A(D),B(D)$ are candidate subnormalized components of orders $A\prec B,B\prec A$, respectively; juxtaposition of a map and an operator means map application. Their sum is $I/4+D$. Each remains in the process linear subspace, has trace two, and obeys $R_BA(D)=A(D)$ or $R_AB(D)=B(D)$, respectively. Norm contraction and the triangle inequality imply
\begin{equation}
 \lambda_{\min}(A(D)),\lambda_{\min}(B(D))
 \ge\frac18-\frac32\|D\|_\infty\ge0.
 \label{eq:CS-ball-positivity}
\end{equation}
Here $\lambda_{\min}$ denotes the smallest eigenvalue. Thus these are positive ordered components, proving the lemma.
\end{proof}

\begin{theorem}[A full-dimensional EC, nonpurifiable neighborhood]
Let $W_0=I/4$ and
\begin{equation}
 \overline W=\frac{21}{22}W_*+\frac1{22}W_0.
 \label{eq:robust-center}
\end{equation}
Here $W_0$ is the uniform process and $\overline W$ the center at noise fraction $1/22$. For every perturbation $\Delta\in\mathcal V_0$,
\begin{equation}
 \|\Delta\|_\infty\le\frac1{264}
 \quad\Longrightarrow\quad
 \overline W+\Delta\in\EC\setminus\CS.
 \label{eq:robust-ball}
\end{equation}
Every process in this ball is nonpurifiable and has no standard deterministic QC-QC realization.
\end{theorem}
\begin{proof}
Set $D=22\Delta$. The lemma gives $W_0+D\in\CS$, and
\begin{equation}
 \overline W+\Delta
 =\frac{21}{22}W_*+\frac1{22}(W_0+D).
 \label{eq:robust-convexity}
\end{equation}
Both normalized components are EC, so convexity proves EC and also ensures positivity and normalization of the perturbed process.

The witness~\eqref{eq:witness} has a positive first term of trace four. Its remaining term is
$(I-\sigma_z^{A_I})\sigma_x^{B_I}\sigma_z^{B_O}/4$, of trace norm four. Hence $\|S\|_1\le8$ by the triangle inequality; its exact value is $2+2\sqrt5$. Using Eq.~\eqref{eq:noise-witness} and trace-norm duality,
\begin{align}
 \Tr[S(\overline W+\Delta)]
 &\le-\frac1{20}+8\|\Delta\|_\infty\notag\\
 &\le-\frac1{20}+\frac8{264}
 =-\frac{13}{660}<0.
 \label{eq:robust-witness}
\end{align}
This excludes CS throughout the closed ball. Equation~\eqref{eq:purifiable-CS} and the QC-QC implication above prove the remaining exclusions.
\end{proof}

More generally, for $0<\eta<1/11$ the same argument proves $W_\eta+\Delta\in\EC\setminus\CS$ whenever
\begin{equation}
 \Delta\in\mathcal V_0,\qquad
 \|\Delta\|_\infty<\min\left\{\frac\eta{12},\frac{1-11\eta}{80}\right\}.
 \label{eq:general-robust-radius}
\end{equation}
These are conservative sufficient radii, not exact boundaries. They establish a nonzero-volume region relative to the trace-fixed valid-process affine space; they do not establish that $W_*$ itself is interior to EC.

For a parameter-space view, set
\begin{equation}
 W(\delta,\alpha,\beta)=\frac14\left[I+\delta\sigma_z^{A_I}
 +\alpha\sigma_z^{A_I}\sigma_x^{B_I}\sigma_z^{B_O}
 +\beta\sum_{k=x,y,z}\sigma_k^{A_O}\sigma_k^{B_I}\right].
 \label{eq:parameter-family}
\end{equation}
Here $\delta,\alpha,\beta$ are varied real coefficients of the same Pauli terms.

The center parameters are $(21/110,63/220,21/110)$. Write $d\delta,d\alpha,d\beta$ for finite offsets from this center, rather than infinitesimal differentials. Since the Pauli sum has norm three, the neighborhood includes all such offsets satisfying
\begin{equation}
 |d\delta|+|d\alpha|+3|d\beta|\le\frac1{66}.
 \label{eq:parameter-neighborhood}
\end{equation}
The full ball also includes valid directions outside this three-parameter family, including ones that break the original flag symmetries.

The conclusions concern finite original qubit ports and arbitrary finite shared-input ancillary assistance. Trusted-input and network certification use additional assumptions~\cite{Trusted,Network}. No claim is made about all certification scenarios, infinite-dimensional processes, or multiple-use composition; a tensor product of mixed-order processes is not automatically a valid process for grouped laboratories.

\section{Reproduction and independent checks}
\emph{Data availability.---}The exact certificates and portable verification code are publicly available in the \href{https://github.com/phycxli/Extensible-Causality-Does-Not-Guarantee-Pure-Quantum-Dilations}{GitHub repository} and in the accompanying supplemental files. They check the process, witness, laboratory responses, and stability without optimization.

The package supplies the original eleven-block certificate in
\path{results/exact_double_flag_hall_sos.json}, its symmetry average in
\path{results/exact_double_flag_hall_sos_symmetric.json}, and the integer tables, ideal multipliers, and principal determinants in the three corresponding \texttt{hall\_sos\_*} JSON files. The three TeX fragments used above are generated directly from those rational data. Exporting an unchanged certificate preserves its frozen bytes and metadata; changed data require a new output directory.

From the project root, the following commands verify both small certificates using an implementation based on Python's standard-library \texttt{Fraction}. They do not import the optimizer or the science generator:
\begingroup\small
\begin{verbatim}
python -m repro.independent_general_ec_certificate_audit \
  --certificate results/exact_double_flag_hall_sos.json \
  --output results/independent_general_ec_certificate_audit.json
python -m repro.independent_general_ec_certificate_audit \
  --certificate results/exact_double_flag_hall_sos_symmetric.json \
  --output results/independent_symmetric_hall_sos_audit.json
python -m repro.independent_counterexample_process_audit
python -m repro.independent_ec_robustness_audit
\end{verbatim}
\endgroup
The continuation syntax shown is for a POSIX shell; the same arguments can be supplied on a single line on Windows. Table regeneration and the independent conic audit are:
\begingroup\small
\begin{verbatim}
python -m repro.export_hall_sos_tables
python results/literature_scope/independent_two_certificate_ec_audit.py
python results/literature_scope/independent_ec_physical_coverage_audit.py
\end{verbatim}
\endgroup
Regeneration uses the supplied project environment with SymPy, NumPy, SciPy, and CVXPY. It performs only exact algebra on the saved certificate, not an optimization. Each exact audit records certificate and verifier SHA256 hashes, equation residuals, and positive rational pivots or determinants.

Independent finite complex mixed-state controls compare the original Born table with Eq.~\eqref{eq:kernel}, perform flag refinement, enumerate all subset cuts, and solve the transportation problem. They include several Kraus branches per report and unequal ancillary dimensions. A deliberately inserted extra transpose fails these controls. The checked examples have Born-reduction residuals below $5\times10^{-17}$ and transportation residuals below $2\times10^{-15}$. Source-hashed computing-node replication runs the same exact certificates; its role is reproducibility, while Eqs.~\eqref{eq:SOS-identity} and \eqref{eq:flow-cuts} establish the universal claim.

The minimal distribution has a single entry point, \texttt{verify.py}, requiring only Python and SymPy. It checks both exact small SOS certificates, physicality and the witness, the ordered split on every valid Pauli basis direction, and the neighborhood constants. It does not invoke the optimizer or regenerate a certificate. The accompanying proof remains necessary for physical coverage, norm contraction, and the externally cited pure-process theorem. A separate review packet lists these obligations for professional mathematical review; such review is not represented as completed by machine verification.

\begin{thebibliography}{99}
\bibitem{Stinespring} W. F. Stinespring, Positive functions on $C^*$-algebras, \href{https://doi.org/10.1090/S0002-9939-1955-0069403-4}{Proc. Am. Math. Soc. \textbf{6}, 211 (1955)}.
\bibitem{QuantumSwitch} G. Chiribella, G. M. D'Ariano, P. Perinotti, and B. Valiron, Quantum computations without definite causal structure, \href{https://doi.org/10.1103/PhysRevA.88.022318}{Phys. Rev. A \textbf{88}, 022318 (2013)}.
\bibitem{QuantumCausality} C. Brukner, Quantum causality, \href{https://doi.org/10.1038/nphys2930}{Nat. Phys. \textbf{10}, 259 (2014)}.
\bibitem{Discriminate} G. Chiribella, Perfect discrimination of no-signalling channels via quantum superposition of causal structures, \href{https://doi.org/10.1103/PhysRevA.86.040301}{Phys. Rev. A \textbf{86}, 040301(R) (2012)}.
\bibitem{GateOrder} M. Ara\'{u}jo, F. Costa, and C. Brukner, Computational Advantage from Quantum-Controlled Ordering of Gates, \href{https://doi.org/10.1103/PhysRevLett.113.250402}{Phys. Rev. Lett. \textbf{113}, 250402 (2014)}.
\bibitem{Communication} P. A. Gu\'{e}rin, A. Feix, M. Ara\'{u}jo, and C. Brukner, Exponential Communication Complexity Advantage from Quantum Superposition of the Direction of Communication, \href{https://doi.org/10.1103/PhysRevLett.117.100502}{Phys. Rev. Lett. \textbf{117}, 100502 (2016)}.
\bibitem{NoisyCommunication} D. Ebler, S. Salek, and G. Chiribella, Enhanced Communication with the Assistance of Indefinite Causal Order, \href{https://doi.org/10.1103/PhysRevLett.120.120502}{Phys. Rev. Lett. \textbf{120}, 120502 (2018)}.
\bibitem{PerfectCommunication} G. Chiribella \emph{et al.}, Indefinite causal order enables perfect quantum communication with zero capacity channels, \href{https://doi.org/10.1088/1367-2630/abe7a0}{New J. Phys. \textbf{23}, 033039 (2021)}.
\bibitem{Procopio} L. M. Procopio \emph{et al.}, Experimental superposition of orders of quantum gates, \href{https://doi.org/10.1038/ncomms8913}{Nat. Commun. \textbf{6}, 7913 (2015)}.
\bibitem{Rubino} G. Rubino \emph{et al.}, Experimental verification of an indefinite causal order, \href{https://doi.org/10.1126/sciadv.1602589}{Sci. Adv. \textbf{3}, e1602589 (2017)}.
\bibitem{Goswami} K. Goswami \emph{et al.}, Indefinite Causal Order in a Quantum Switch, \href{https://doi.org/10.1103/PhysRevLett.121.090503}{Phys. Rev. Lett. \textbf{121}, 090503 (2018)}.
\bibitem{WeiExperiment} K. Wei \emph{et al.}, Experimental Quantum Switching for Exponentially Superior Quantum Communication Complexity, \href{https://doi.org/10.1103/PhysRevLett.122.120504}{Phys. Rev. Lett. \textbf{122}, 120504 (2019)}.
\bibitem{Supermaps} G. Chiribella, G. M. D'Ariano, and P. Perinotti, Transforming quantum operations: Quantum supermaps, \href{https://doi.org/10.1209/0295-5075/83/30004}{Europhys. Lett. \textbf{83}, 30004 (2008)}.
\bibitem{CombArchitecture} G. Chiribella, G. M. D'Ariano, and P. Perinotti, Quantum Circuit Architecture, \href{https://doi.org/10.1103/PhysRevLett.101.060401}{Phys. Rev. Lett. \textbf{101}, 060401 (2008)}.
\bibitem{QuantumNetworks} G. Chiribella, G. M. D'Ariano, and P. Perinotti, Theoretical framework for quantum networks, \href{https://doi.org/10.1103/PhysRevA.80.022339}{Phys. Rev. A \textbf{80}, 022339 (2009)}.
\bibitem{OCB} O. Oreshkov, F. Costa, and C. Brukner, Quantum correlations with no causal order, \href{https://doi.org/10.1038/ncomms2076}{Nat. Commun. \textbf{3}, 1092 (2012)}.
\bibitem{PurificationAxiom} G. Chiribella, G. M. D'Ariano, and P. Perinotti, Probabilistic theories with purification, \href{https://doi.org/10.1103/PhysRevA.81.062348}{Phys. Rev. A \textbf{81}, 062348 (2010)}.
\bibitem{purif} M. Ara\'{u}jo, A. Feix, M. Navascu\'{e}s, and C. Brukner, A purification postulate for quantum mechanics with indefinite causal order, \href{https://doi.org/10.22331/q-2017-04-26-10}{Quantum \textbf{1}, 10 (2017)}.
\bibitem{BranciardCI} C. Branciard, M. Ara\'{u}jo, A. Feix, F. Costa, and C. Brukner, The simplest causal inequalities and their violation, \href{https://doi.org/10.1088/1367-2630/18/1/013008}{New J. Phys. \textbf{18}, 013008 (2016)}.
\bibitem{Witness} M. Ara\'{u}jo \emph{et al.}, Witnessing causal nonseparability, \href{https://doi.org/10.1088/1367-2630/17/10/102001}{New J. Phys. \textbf{17}, 102001 (2015)}.
\bibitem{WitnessGuide} C. Branciard, Witnesses of causal nonseparability: an introduction and a few case studies, \href{https://doi.org/10.1038/srep26018}{Sci. Rep. \textbf{6}, 26018 (2016)}.
\bibitem{Feix} A. Feix, M. Ara\'{u}jo, and C. Brukner, Causally nonseparable processes admitting a causal model, \href{https://doi.org/10.1088/1367-2630/18/8/083040}{New J. Phys. \textbf{18}, 083040 (2016)}.
\bibitem{EC} O. Oreshkov and C. Giarmatzi, Causal and causally separable processes, \href{https://doi.org/10.1088/1367-2630/18/9/093020}{New J. Phys. \textbf{18}, 093020 (2016)}.
\bibitem{Yokojima} W. Yokojima, M. T. Quintino, A. Soeda, and M. Murao, Consequences of preserving reversibility in quantum superchannels, \href{https://doi.org/10.22331/q-2021-04-26-441}{Quantum \textbf{5}, 441 (2021)}.
\bibitem{Cyclic} J. Barrett, R. Lorenz, and O. Oreshkov, Cyclic quantum causal models, \href{https://doi.org/10.1038/s41467-020-20456-x}{Nat. Commun. \textbf{12}, 885 (2021)}.
\bibitem{TimeDelocalized} O. Oreshkov, Time-delocalized quantum subsystems and operations: on the existence of processes with indefinite causal structure in quantum mechanics, \href{https://doi.org/10.22331/q-2019-12-02-206}{Quantum \textbf{3}, 206 (2019)}.
\bibitem{Wechs} J. Wechs, A. A. Abbott, and C. Branciard, On the definition and characterisation of multipartite causal (non)separability, \href{https://doi.org/10.1088/1367-2630/aaf352}{New J. Phys. \textbf{21}, 013027 (2019)}.
\bibitem{Wei} W. Wei and S. Pang, Extensibly Causally Separable Processes Admit Realizations as Quantum Circuits with Classical Control of Causal Order, arXiv:2609.18559v1 (2026).
\bibitem{Wechs2026} J. Wechs, A. A. Abbott, and C. Branciard, All causally separable quantum processes are quantum circuits with classical control of causal order, arXiv:2609.20774v1 (2026).
\bibitem{SM} See Supplemental Material for process physicality, the general Born reduction, the transportation theorem, all rational SOS matrices, and independent verification commands.
\bibitem{Trusted} H. Dourdent \emph{et al.}, Semi-Device-Independent Certification of Causal Nonseparability with Trusted Quantum Inputs, \href{https://doi.org/10.1103/PhysRevLett.129.090402}{Phys. Rev. Lett. \textbf{129}, 090402 (2022)}.
\bibitem{Network} H. Dourdent \emph{et al.}, Network-Device-Independent Certification of Causal Nonseparability, \href{https://doi.org/10.22331/q-2024-10-30-1514}{Quantum \textbf{8}, 1514 (2024)}.
\bibitem{QCQC} J. Wechs, H. Dourdent, A. A. Abbott, and C. Branciard, Quantum circuits with classical versus quantum control of causal order, \href{https://doi.org/10.1103/PRXQuantum.2.030335}{PRX Quantum \textbf{2}, 030335 (2021)}.
\end{thebibliography}
\end{document}
